\documentclass[11pt]{article}
\usepackage[margin=1in]{geometry}
\usepackage{amsmath,amssymb,booktabs,hyperref,xcolor,listings}
\lstset{basicstyle=\ttfamily\small,breaklines=true,columns=fullflexible}
\title{PipelineHalt: Halting on Paper--Lean Drift and Asking Instead of Guessing}
\author{Benjamin Stanley Frohman\\\texttt{frohmanbenjamin@gmail.com}\\ORCID: \href{https://orcid.org/0009-0006-7068-3718}{0009-0006-7068-3718}}
\date{November 2026}
\begin{document}
\maketitle
\begin{abstract}
A formalization is only as faithful as the match between the statements it proves and the claims the accompanying paper makes. These two artifacts drift: indices shift, hypotheses are added, sets are re-encoded, numbering diverges. We describe PipelineHalt, a small open-source tool that treats the paper's claims as an anchor and the Lean statements as the live artifact, flags places where the Lean is \emph{silent}, \emph{narrower}, or \emph{conditional} relative to the paper, or where the paper itself is \emph{ambiguous}, and then halts with a precise clarification question rather than correcting anything. Answers are recorded verbatim in the anchor under stable identifiers so questions are not re-asked. The tool does not and cannot recover intent that was never written down; its matching is heuristic and syntactic. We report its real output on a small Lean development, including a false positive.
\end{abstract}

\section{Introduction}
Machine-checked proofs in Lean~4~\cite{moura2021lean4} with Mathlib~\cite{mathlib2020} certify that a \emph{Lean statement} holds. Readers, however, consume the \emph{paper}. Whether the paper's Theorem~3.1 and the Lean \texttt{two\_le\_gap} say the same thing is a question the kernel does not ask. Large-language-model agents that write or revise either artifact make the problem sharper: they will happily ``fix'' a mismatch in whichever direction is locally convenient.

PipelineHalt is an extension of RecurringAlignment~\cite{frohman2026ra}. Its design principle is simple: \textbf{on a detected difference, stop and ask the author}. It never edits a claim or a Lean file.

\section{The residual-gap problem}
Call the \emph{residual gap} the set of differences between what the paper asserts and what the Lean proves that survive compilation. Some are benign conventions (0- vs.\ 1-indexed primes); some change the mathematics (an extra hypothesis). Syntactically the two cases often look identical, and distinguishing them requires information that lives only in the author's head. A tool can therefore at best \emph{surface} the gap and \emph{elicit} a decision; it cannot close it alone. We make that limitation the interface.

\section{Design}
\paragraph{Inputs.} The anchor is a list of claims (id, printed number, title, statement, optional explicit \texttt{lean:} mapping and stated assumptions), read from JSON, Markdown, or LaTeX theorem environments (with amsthm section numbering reproduced); PDF is supported best-effort via \texttt{pdftotext}. Lean files are parsed with a lightweight bracket-aware scanner extracting declarations, binders, type, docstring, \texttt{sorry}, and use of project axioms.

\paragraph{Mapping file.} The practical way to cut first-run noise is a small, hand-written \emph{mapping file}, \texttt{pipelinehalt.map.json}, kept separate from the extracted claims. It maps each claim id to Lean declaration names, optionally lists \emph{aliases} (e.g.\ a docstring's ``Theorem~1'' for the paper's Theorem~3.1), and may list declarations deliberately left unmapped. It is found by default beside the anchor or in a parent directory up to the repository root, and can be overridden (\texttt{--map}) or disabled (\texttt{--no-map}). Mapped claims skip heuristic pairing, and aliases suppress numbering questions; all other checks still run on the mapped declarations. The mapping is human-authored: it says \emph{where to look}, reduces pairing noise, and is not validated semantically. A wrong mapping yields wrong findings.

\paragraph{Matching.} Without a mapping, each declaration is assigned to the claim maximizing
$0.6\,D(\text{sig}) + 0.4\,D(\text{doc}) + 0.3\,[\text{doc cites the claim's number}]$,
where $D$ is the Dice coefficient~\cite{dice1945} over stemmed tokens; definitions match definition-like claims only. This is a heuristic and is the main source of error.

\paragraph{Checks.} Per claim:
\begin{itemize}
\item \textbf{Silent}: a mapped name is missing, or nothing matches.
\item \textbf{Narrower}: a Lean lower bound on a variable differs from the claim's, or a set literal in the claim appears in no matched statement.
\item \textbf{Conditional}: a hypothesis binder or premise not covered by the claim text or its stated assumptions; \texttt{sorry}; a project axiom.
\item \textbf{Ambiguous}: hedging words, empty statement, title unrelated to statement, duplicate ids, or a docstring theorem number disagreeing with the printed one.
\end{itemize}
Each mismatch has an id $=\mathrm{sha256}(\text{claim},\text{kind},\text{decl},\text{key})$ truncated to 12 hex digits, and a weight (silent 1, narrower 0.6, conditional 0.5, sorry 0.8, axiom 0.7, ambiguous 0.3); claim drift is the capped sum.

\section{Clarification protocol}
Unanswered mismatches yield status \textsc{halt} (exit code 2). Interactively, questions are asked on the terminal; non-interactively, they are written to a pending file to be filled in and imported. An answer is stored verbatim with a verdict parsed from its first word: \emph{intentional} (weight 0), \emph{unintended} (full weight), or other (half weight). Since ids are stable, an answered question is never re-asked; if the Lean changes so that the key changes, a new question appears. For agents, a JSON protocol (\texttt{pipelinehalt/1}) returns the questions and accepts answers; the accompanying prompt template instructs agents to relay questions unchanged and never answer them. A CI mode emits GitHub annotations and fails on halts or drift above a threshold.

\section{Example: DHL[3,2] and prime gaps}
We ran PipelineHalt 0.4.0 on the repository for \cite{frohman2026dhl} (commit \texttt{69a13d3}): 9 claims extracted verbatim from \texttt{paper/DHL32\_note.tex}, 34 Lean declarations. Table~\ref{tab:runs} compares a run without a mapping with a run using the mapping file shipped in \texttt{examples/dhl32/}, which we drafted from the Lean docstrings.
\begin{table}[h]\centering\small
\begin{tabular}{lcccc}\toprule
Run & Status & Drift & Open questions & Unmatched-citation notes \\\midrule
No mapping (\texttt{--no-map}) & \textsc{halt} & 0.267 & 5 & 3 \\
Mapping file (default) & \textsc{halt} & 0.160 & 3 & 0 \\\bottomrule
\end{tabular}
\caption{Real output on the DHL[3,2] development.}\label{tab:runs}
\end{table}

\paragraph{Without a mapping.} The heuristic paired \texttt{Ptoy\_DHL32} with Theorem~3.2. This produced a \emph{false} ``silent'' finding for Theorem~3.3 and a spurious unstated-hypothesis question (\texttt{H.card = 3}) from a wrongly paired declaration. It also asked three docstring-numbering questions (``Theorem 1/2/4''). It \emph{missed} the $n\ge1$ vs.\ $n\ge2$ difference, because \texttt{two\_le\_gap} was not paired with Theorem~3.1.

\paragraph{With the mapping file.} Three questions remain:
\begin{itemize}
\item Theorem~3.1: \texttt{two\_le\_gap} and \texttt{gap\_even} assume $1\le n$, while the paper says $n\ge2$ (two questions);
\item Theorem~3.3: the paper's $\{0,1\}$ appears in no mapped Lean statement.
\end{itemize}
The second is very likely a false positive: \texttt{Ptoy} is defined by \texttt{n \% 4 < 2}, which is equivalent to $n \bmod 4\in\{0,1\}$ over $\mathbb{N}$, but a syntactic check cannot establish that. The first may be an indexing convention; only the author can say. The mapping removed pairing and numbering noise; it did not, and cannot, resolve these. They remain unanswered in the published example; we did not answer them on the author's behalf.

\section{Limitations}
PipelineHalt is syntactic. It does not understand mathematics, does not check proofs, and does not recover unstated intent; an aligned report means only that these checks found nothing. Token matching mis-assigns declarations unless a mapping file is supplied, and the mapping itself is unverified; bound and set checks miss equivalent encodings and miss differences expressed otherwise; LaTeX macro normalization is partial; PDF extraction is fragile. Drift weights are arbitrary conventions, not measurements. The evaluation is a single small example.

\section{Conclusion}
PipelineHalt turns paper--Lean drift into explicit, recorded questions with stable ids, and refuses to resolve them itself. The value lies not in the detection heuristics, which are modest, but in the discipline: differences are surfaced, the pipeline halts, and the author's answer becomes part of the record.

\paragraph{Availability.} Source, tests and example outputs: \url{https://github.com/BenFrohman/PipelineHalt} (code MIT; paper CC BY 4.0), derived from \texttt{packages/pipelinehalt} in \cite{frohman2026ra}.

\bibliographystyle{plain}
\bibliography{pipelinehalt}
\end{document}
